\begin{proof}[Proof of \cref{obj:thm:full-design-lower}]
Fix \(s\in(0,1]\) and integers \(n,d\geq2\). We use the parameters and features of \cref{obj:def:cosine-prior}, written explicitly as
\[
B=\binom d2,\qquad
L=\max\left\{1,\left\lceil\sqrt{4096n/B}\right\rceil\right\},
\qquad
M=BL^2,\qquad
C_0=\sqrt{48+32\pi^2},
\]
\[
\mathcal I_{d,L}
=\{(j,h,k,\ell):1\leq j<h\leq d,\ 1\leq k,\ell\leq L\},
\qquad
V_\alpha(x)=2\cos(\pi kx_j)\cos(\pi\ell x_h),
\qquad
V(x)=(V_\alpha(x))_{\alpha\in\mathcal I_{d,L}}.
\]
Thus \(B>0\), \(L\geq1\), \(M>0\), and \(C_0>0\).

\begin{enumerate}
\item
The constant
\[
c_s=\frac{2}{(48+32\pi^2)\,16384^s}>0
\]
depends on \(s\) alone. We first compare it with the cutoff normalization. Define
\[
\eta_{\mathrm{cut}}=\frac{4096n}{B},
\qquad
\Gamma_{\mathrm{cut}}=\max\{1,n/B\},
\qquad
\Lambda_{\mathrm{cut}}=16384^s.
\]
Since \(\eta_{\mathrm{cut}}\leq4096\Gamma_{\mathrm{cut}}\) and
\(\Gamma_{\mathrm{cut}}\geq1\), the ceiling bound gives
\[
\sqrt{\eta_{\mathrm{cut}}}\leq64\sqrt{\Gamma_{\mathrm{cut}}},
\qquad
\left\lceil\sqrt{\eta_{\mathrm{cut}}}\right\rceil
\leq\sqrt{\eta_{\mathrm{cut}}}+1,
\]
and hence
\[
L\leq64\sqrt{\Gamma_{\mathrm{cut}}}+1
\leq128\sqrt{\Gamma_{\mathrm{cut}}},
\qquad
L^2\leq16384\Gamma_{\mathrm{cut}}.
\]
Raising the last inequality to the positive power \(s\) yields
\[
L^{2s}=(L^2)^s
\leq\Lambda_{\mathrm{cut}}\Gamma_{\mathrm{cut}}^s.
\]

We next verify the scale comparison in its two cases. If \(n/B\leq1\), then
\[
b_s(n,d)\Gamma_{\mathrm{cut}}^s=b_s(n,d)\leq1.
\]
If \(n/B>1\), then
\[
\begin{aligned}
b_s(n,d)\Gamma_{\mathrm{cut}}^s
&\leq (B/n)^s(n/B)^s\\
&=\bigl((B/n)(n/B)\bigr)^s=1.
\end{aligned}
\]
Because \(b_s(n,d)\geq0\) and \(\Lambda_{\mathrm{cut}}>0\), these inequalities imply
\[
L^{2s}b_s(n,d)
\leq\Lambda_{\mathrm{cut}}
       b_s(n,d)\Gamma_{\mathrm{cut}}^s
\leq\Lambda_{\mathrm{cut}}.
\]
Using \(C_0^2=48+32\pi^2\) and positive denominators, we obtain
\[
c_sb_s(n,d)
=\frac{2b_s(n,d)}{C_0^2\Lambda_{\mathrm{cut}}}
\leq\frac{2}{C_0^2L^{2s}}.
\]
% lean: CausalSmith.Experimentation.BivariateSobolevDesignCapacity.priorCutoff_sq_upper
% lean: CausalSmith.Experimentation.BivariateSobolevDesignCapacity.priorCutoff_normalized_lower

\item
Fix an admissible design \(\pi\in\mathcal D_{n,d,s}\). With \(X\) drawn as in
\cref{obj:ass:uniform-draw} and \(Z\) drawn conditionally from \(\pi(X)\), write the scaled signing loss appearing in \cref{obj:def:criterion} as
\[
\mathcal L_{n,d,s}(\pi,m)
=\frac4n\mathbb E_{X,Z}
 \left[\left(\sum_{i=1}^n Z_i m(X_i)\right)^2\right].
\]
The coefficient signs have the independent uniform law of
\cref{obj:def:cosine-prior}, independently of this experiment. Their second moments satisfy
\[
\mathbb E_\xi[\xi_\alpha\xi_\beta]
=
\begin{cases}
1,&\alpha=\beta,\\
0,&\alpha\ne\beta.
\end{cases}
\]
For the first case, \(\xi_\alpha^2=1\). For the second, flipping the sign at
\(\alpha\) preserves the uniform coefficient law and reverses the product
\(\xi_\alpha\xi_\beta\), so its mean is zero.

For every deterministic sample
\(x=(x_1,\ldots,x_n)\in([0,1]^d)^n\) and signing
\(z\in\{-1,1\}^n\), the prior formula gives
\[
\sum_{i=1}^n z_i m_\xi(x_i)
=\frac{1}{C_0L^s\sqrt M}
 \sum_{\alpha\in\mathcal I_{d,L}}
 \xi_\alpha\sum_{i=1}^n z_iV_\alpha(x_i).
\]
Expanding the square and applying the second-moment identity therefore gives
\[
\begin{aligned}
\mathbb E_\xi
 \left[\left(\sum_{i=1}^n z_i m_\xi(x_i)\right)^2\right]
&=\frac{1}{C_0^2L^{2s}M}
 \sum_{\alpha,\beta\in\mathcal I_{d,L}}
 \mathbb E_\xi[\xi_\alpha\xi_\beta]
 \left(\sum_i z_iV_\alpha(x_i)\right)
 \left(\sum_i z_iV_\beta(x_i)\right)\\
&=\frac{1}{C_0^2L^{2s}M}
 \sum_{\alpha\in\mathcal I_{d,L}}
 \left(\sum_i z_iV_\alpha(x_i)\right)^2\\
&=\frac{1}{C_0^2L^{2s}M}
 \left\|\sum_i z_iV(x_i)\right\|_2^2.
\end{aligned}
\]
Finite coefficient averaging commutes with the nonnegative sample and assignment integrals. Consequently,
\[
\mathbb E_\xi\mathcal L_{n,d,s}(\pi,m_\xi)
=\frac{4}{nC_0^2L^{2s}M}
 \mathbb E_{X,Z}
 \left\|\sum_i Z_iV(X_i)\right\|_2^2.
\]
For each fixed sample \(x\), the probability kernel satisfies
\[
\int_{\{-1,1\}^n}
 \left\|\sum_i z_iV(x_i)\right\|_2^2\,\pi(x,dz)
\geq
\min_{z\in\{-1,1\}^n}
 \left\|\sum_i z_iV(x_i)\right\|_2^2.
\]
Integrating this pointwise comparison yields
\[
\mathbb E_\xi\mathcal L_{n,d,s}(\pi,m_\xi)
\geq
\frac{4}{nC_0^2L^{2s}M}
\mathbb E_X
 \min_{z\in\{-1,1\}^n}
 \left\|\sum_i z_iV(X_i)\right\|_2^2.
\]
% lean: CausalSmith.Experimentation.BivariateSobolevDesignCapacity.coefficientSigns_product_integral
% lean: CausalSmith.Experimentation.BivariateSobolevDesignCapacity.coefficientSigns_sum_second_moment
% lean: CausalSmith.Experimentation.BivariateSobolevDesignCapacity.mXi_signed_sum_prior_second_moment
% lean: CausalSmith.Experimentation.BivariateSobolevDesignCapacity.priorRisk_eq_feature_energy
% lean: CausalSmith.Experimentation.BivariateSobolevDesignCapacity.priorRisk_ge_min_feature_energy

\item
The cutoff also supplies the feature-dimension threshold. Indeed,
\[
\sqrt{4096n/B}
\leq\left\lceil\sqrt{4096n/B}\right\rceil
\leq L.
\]
Squaring these nonnegative quantities and multiplying by \(B>0\) gives
\[
\frac{4096n}{B}\leq L^2,
\qquad
4096n\leq BL^2=M.
\]
Under \cref{obj:ass:uniform-draw}, the feature rows \(V(X_i)\) are independent copies. Their coordinates are measurable and square-integrable, and
\cref{obj:lem:cosine-prior} supplies
\[
\mathbb E\!\left[V(X_1)V(X_1)^{\mathsf T}\right]=I_M.
\]
The expected absolute-supremum contraction inequality of
\citep[Theorem~12(4), with Rademacher complexity as in Definition~2]{BartlettMendelson2002Complexities}
supplies the contraction step underlying \cref{obj:lem:isotropic-row-signing}, for its finite classes and finite-dimensional Euclidean unit-ball linear class. Applying that signing bound with feature dimension \(M\geq4096n\) gives
\[
\mathbb E_X
 \min_{z\in\{-1,1\}^n}
 \left\|\sum_i z_iV(X_i)\right\|_2^2
\geq\frac{nM}{2}.
\]
The finite minimum is measurable, so this is the same sample integral that appears in the preceding step.
% lean: CausalSmith.Experimentation.BivariateSobolevDesignCapacity.priorCutoff_dimension_lower
% lean: CausalSmith.Experimentation.BivariateSobolevDesignCapacity.prior_min_feature_energy_lower

\item
Substituting the feature-energy bound into the prior-risk comparison and cancelling the positive factors gives
\[
\begin{aligned}
\mathbb E_\xi\mathcal L_{n,d,s}(\pi,m_\xi)
&\geq
\frac4n\left(\frac{1}{C_0L^s\sqrt M}\right)^2
\frac{nM}{2}\\
&=\frac{2}{C_0^2L^{2s}}\\
&\geq c_sb_s(n,d).
\end{aligned}
\]
The final inequality uses the cutoff comparison \(c_sb_s(n,d)\leq 2/(C_0^2L^{2s})\). This proves the asserted prior lower bound for every admissible design.
% lean: CausalSmith.Experimentation.BivariateSobolevDesignCapacity.priorRisk_ge_normalized_cutoff
% lean: CausalSmith.Experimentation.BivariateSobolevDesignCapacity.full_design_lower

\item
Since \(d\geq2\), \(L\geq1\), and \(0<s\leq1\),
\cref{obj:lem:cosine-prior} ensures that every prior realization belongs to
\(\mathcal H_{d,s}\). The coefficient index set has cardinality \(M\), so its uniform law has total mass
\[
\sum_{\xi\in\{-1,1\}^{\mathcal I_{d,L}}}2^{-M}=1.
\]
For each fixed admissible \(\pi\), every prior atom satisfies
\[
\mathcal L_{n,d,s}(\pi,m_\xi)
\leq
\sup_{m\in\mathcal H_{d,s}}\mathcal L_{n,d,s}(\pi,m).
\]
Multiplying by the prior weights and summing therefore yields
\[
\begin{aligned}
\mathbb E_\xi\mathcal L_{n,d,s}(\pi,m_\xi)
&=\sum_{\xi\in\{-1,1\}^{\mathcal I_{d,L}}}
 2^{-M}\mathcal L_{n,d,s}(\pi,m_\xi)\\
&\leq
\left(\sum_{\xi\in\{-1,1\}^{\mathcal I_{d,L}}}2^{-M}\right)
\sup_{m\in\mathcal H_{d,s}}\mathcal L_{n,d,s}(\pi,m)\\
&=\sup_{m\in\mathcal H_{d,s}}\mathcal L_{n,d,s}(\pi,m).
\end{aligned}
\]
Thus every admissible design has worst-case loss at least \(c_sb_s(n,d)\). Taking the infimum in \cref{obj:def:criterion} gives
\[
R_s(n,d)
=\inf_{\pi\in\mathcal D_{n,d,s}}
 \sup_{m\in\mathcal H_{d,s}}\mathcal L_{n,d,s}(\pi,m)
\geq c_sb_s(n,d).
\]
% lean: CausalSmith.Experimentation.BivariateSobolevDesignCapacity.priorRisk_le_worstLoss_of_legal
% lean: CausalSmith.Experimentation.BivariateSobolevDesignCapacity.capacity_lower_of_priorRisk
\end{enumerate}
\end{proof}
